\chapter{Capitolo 6 - Funtori, naturali, aggiunzioni, monadi/comonadi}

\section{1. Problema}

Dopo aver esteso in chiave ordinativa oggetti, morfismi, composizione, identità, isomorfismi, equivalenze, limiti e colimiti, resta da affrontare il livello trasformazionale che rende davvero potente la category theory:

\begin{enumerate}
\def\labelenumi{\arabic{enumi}.}
\tightlist
\item
  funtori;
\item
  trasformazioni naturali;
\item
  aggiunzioni;
\item
  monadi e comonadi.
\end{enumerate}

Nel quadro classico, questi costrutti sono il motore del trasferimento di struttura tra domini diversi. In OCT dobbiamo rispondere a una domanda più severa: non solo ``il trasferimento è formalmente corretto?'', ma ``il trasferimento preserva o degrada la validità ordinativa nel contesto osservativo?''.

Il problema del capitolo è quindi:

\textbf{come estendere il livello trasformazionale classico in modo che il trasporto di struttura non cancelli informazione ordinativa critica né introduca falsa equivalenza operativa?}

La difficoltà è strutturale:

\begin{enumerate}
\def\labelenumi{\arabic{enumi}.}
\tightlist
\item
  un funtore può preservare composizione e identità classiche ma deteriorare coerenza o emergenza;
\item
  una naturalità formalmente impeccabile può essere ordinativamente sterile;
\item
  un'aggiunzione può apparire elegante ma non robusta su domini reali;
\item
  una monade può stabilizzare il flusso sintattico e contemporaneamente amplificare degenerazione funzionale.
\end{enumerate}

Per questo il capitolo adotta una regola metodologica unica:

\begin{enumerate}
\def\labelenumi{\arabic{enumi}.}
\tightlist
\item
  conservazione classica come prerequisito;
\item
  preservazione ordinativa come condizione d'uso reale.
\end{enumerate}

\subsection{1.1 Perché il livello trasformazionale è decisivo}

Se OCT restasse confinata a proprietà interne di una singola categoria, il suo impatto sarebbe limitato. La scienza contemporanea vive invece di traduzioni:

\begin{enumerate}
\def\labelenumi{\arabic{enumi}.}
\tightlist
\item
  da teoria a implementazione;
\item
  da dominio a dominio;
\item
  da rappresentazione a rappresentazione.
\end{enumerate}

Il livello funtoriale e naturale è il luogo in cui queste traduzioni avvengono. Se il capitolo non chiarisce come misurare la qualità ordinativa del trasporto, il rischio è alto:

\begin{enumerate}
\def\labelenumi{\arabic{enumi}.}
\tightlist
\item
  formalismi impeccabili con applicazioni fragili;
\item
  applicazioni convincenti senza fondazione trasferibile.
\end{enumerate}

\subsection{1.2 Obiettivo tecnico del capitolo}

L'obiettivo non è riscrivere tutta la teoria classica delle trasformazioni, ma costruire un layer ordinativo minimo e riusabile che specifichi:

\begin{enumerate}
\def\labelenumi{\arabic{enumi}.}
\tightlist
\item
  quando un funtore è ordinativamente ammissibile;
\item
  quando una naturalezza è anche ordinativamente significativa;
\item
  quando una aggiunzione è operativamente reale;
\item
  quando una monade/comonade produce accumulo costruttivo e non deriva degenerativa.
\end{enumerate}

\begin{center}\rule{0.5\linewidth}{0.5pt}\end{center}

\section{2. Tesi del capitolo}

La tesi è articolata in dieci enunciati.

\begin{enumerate}
\def\labelenumi{\arabic{enumi}.}
\tightlist
\item
  Un funtore ordinativo preserva struttura classica e mantiene condizioni minime su \texttt{Coh/Phi}.
\item
  Le trasformazioni naturali devono essere valutate anche per coerenza ordinativa del quadrato naturale.
\item
  La naturalità classica non implica automaticamente naturalità ordinativa.
\item
  Le aggiunzioni ordinative richiedono controllo del trasporto in entrambe le direzioni.
\item
  Monadi e comonadi ordinative sono valide solo se l'iterazione non produce deriva sistematica verso sterilità o degenerazione.
\item
  La distinzione tra preservazione formale e preservazione ordinativa è non ridondante.
\item
  Esiste una nozione utile di ``funtore dimenticante ordinativamente sicuro''.
\item
  Il criterio locale sui singoli morfismi non basta: serve controllo di traiettorie trasformazionali composte.
\item
  Le nozioni introdotte sono compatibili con il recupero classico.
\item
  Il capitolo prepara direttamente il Capitolo 7 su monoidali/enriched/fibrations/topos.
\end{enumerate}

\begin{center}\rule{0.5\linewidth}{0.5pt}\end{center}

\section{3. Definizioni canoniche}

\subsection{Definizione 6.1 (Funtore ordinativamente ammissibile)}

Sia \texttt{F:\ C\_Omega\ -\textgreater{}\ D\_Omega\textquotesingle{}} un funtore classico. Diciamo che \texttt{F} è ordinativamente ammissibile se:

\begin{enumerate}
\def\labelenumi{\arabic{enumi}.}
\tightlist
\item
  preserva composizione e identità classiche;
\item
  per la famiglia di oggetti/morfismi rilevanti, la mappa indotta non viola soglie minime di coerenza;
\item
  non annulla sistematicamente la funzione emergente su classi operative dichiarate.
\end{enumerate}

\subsection{Definizione 6.2 (Indice di preservazione ordinativa funtoriale)}

Definiamo un indice sintetico:

\texttt{Pres\_F\ :=\ \textless{}DeltaCoh\_F,\ DeltaPhi\_F,\ Drift\_F\textgreater{}}

dove:

\begin{enumerate}
\def\labelenumi{\arabic{enumi}.}
\tightlist
\item
  \texttt{DeltaCoh\_F} misura variazione media di coerenza pre/post trasporto;
\item
  \texttt{DeltaPhi\_F} misura variazione emergenziale;
\item
  \texttt{Drift\_F} misura tendenza di lungo periodo a spostare stati verso B/C.
\end{enumerate}

\subsection{Definizione 6.3 (Trasformazione naturale ordinativa)}

Data \texttt{eta:\ F\ =\textgreater{}\ G}, la trasformazione è ordinativamente valida se:

\begin{enumerate}
\def\labelenumi{\arabic{enumi}.}
\tightlist
\item
  è naturale in senso classico;
\item
  i quadrati naturali coinvolti restano in classe A o in banda di robustezza dichiarata;
\item
  non introducono perdita emergenziale critica lungo i componenti di \texttt{eta}.
\end{enumerate}

\subsection{Definizione 6.4 (Aggiunzione ordinativa)}

Una aggiunzione \texttt{F\ ⊣\ G} è ordinativa se:

\begin{enumerate}
\def\labelenumi{\arabic{enumi}.}
\tightlist
\item
  l'aggiunzione classica è valida;
\item
  unità e counità risultano ordinativamente non-degenerative;
\item
  il ciclo avanti-indietro non produce accumulo di perdita oltre soglia.
\end{enumerate}

\subsection{Definizione 6.5 (Monade ordinativa)}

Data monade \texttt{T=(T,eta,mu)}, diciamo che è ordinativa se:

\begin{enumerate}
\def\labelenumi{\arabic{enumi}.}
\tightlist
\item
  i dati monadici classici sono soddisfatti;
\item
  l'iterazione di \texttt{T} non induce deriva sistematica su \texttt{Coh/Phi} oltre limiti dichiarati;
\item
  i punti fissi operativi non collassano in sterilità generalizzata.
\end{enumerate}

\subsection{Definizione 6.6 (Comonade ordinativa)}

Data comonade \texttt{W=(W,epsilon,delta)}, diciamo che è ordinativa se:

\begin{enumerate}
\def\labelenumi{\arabic{enumi}.}
\tightlist
\item
  i dati comonadici classici sono soddisfatti;
\item
  il processo di coestrazione/duplicazione preserva margine ordinativo;
\item
  non introduce rumore strutturale che annulli funzione emergente.
\end{enumerate}

\subsection{Definizione 6.7 (Funtore dimenticante sicuro)}

Un funtore dimenticante è ordinativamente sicuro su classe \texttt{K} se la perdita informativa indotta è preservativa e non cambia classificazione A/B/C oltre una soglia di tolleranza dichiarata.

\subsection{Definizione 6.8 (Catena trasformazionale)}

Una catena trasformazionale è una composizione di funtori e naturali:

\texttt{Xi\ =\ H\_n\ \ o\ \ ...\ \ o\ \ H\_1} con naturali intermedie.

La catena è ordinativamente affidabile se ogni prefisso mantiene margine positivo e robustezza minima.

\begin{center}\rule{0.5\linewidth}{0.5pt}\end{center}

\section{4. Sviluppo formale}

\subsection{4.1 Funtori: preservare forma non basta}

Nel classico un funtore è definito da:

\begin{enumerate}
\def\labelenumi{\arabic{enumi}.}
\tightlist
\item
  mappa su oggetti;
\item
  mappa su morfismi;
\item
  preservazione di composizione e identità.
\end{enumerate}

In OCT questa definizione resta necessaria ma non sufficiente. Due funtori possono essere entrambi validi classicamente e differire drasticamente su qualità ordinativa del trasporto.

Per questo introduciamo il criterio duale:

\begin{enumerate}
\def\labelenumi{\arabic{enumi}.}
\tightlist
\item
  \textbf{validità funtoriale classica} (\texttt{Fun\_class});
\item
  \textbf{validità funtoriale ordinativa} (\texttt{Fun\_ord}).
\end{enumerate}

Solo la congiunzione \texttt{Fun\_class\ \ AND\ \ Fun\_ord} autorizza uso forte del trasporto in inferenze operative.

\subsection{4.2 Naturali: quadrato commutativo vs quadrato fertile}

Una trasformazione naturale classica garantisce commutazione dei quadrati. In OCT aggiungiamo una domanda:

\begin{enumerate}
\def\labelenumi{\arabic{enumi}.}
\tightlist
\item
  i quadrati commutano preservando anche qualità ordinativa?
\end{enumerate}

Possibile anomalia:

\begin{enumerate}
\def\labelenumi{\arabic{enumi}.}
\tightlist
\item
  il quadrato commuta formalmente;
\item
  uno dei percorsi produce perdita emergenziale critica;
\item
  il risultato è ``naturale'' ma non affidabile in senso ordinativo.
\end{enumerate}

Da qui il criterio:

\begin{enumerate}
\def\labelenumi{\arabic{enumi}.}
\tightlist
\item
  naturalità classica necessaria;
\item
  naturalità ordinativa verificata su classi di casi rilevanti.
\end{enumerate}

\subsection{4.3 Aggiunzioni: simmetria utile o ciclo dissipativo}

L'aggiunzione \texttt{F\ ⊣\ G} è spesso letta come equilibrio elegante tra due direzioni di traduzione. In OCT questo equilibrio può essere apparente se:

\begin{enumerate}
\def\labelenumi{\arabic{enumi}.}
\tightlist
\item
  \texttt{F} preserva bene \texttt{Coh} ma degrada \texttt{Phi};
\item
  \texttt{G} recupera forma ma accumula drift su iterazioni.
\end{enumerate}

Per questo valutiamo unità e counità non solo come dati formali, ma come circuiti operativi. Un ciclo di round-trip deve essere controllato per:

\begin{enumerate}
\def\labelenumi{\arabic{enumi}.}
\tightlist
\item
  perdita coerentiva cumulativa;
\item
  perdita emergenziale cumulativa;
\item
  stabilità su vicinato contestuale.
\end{enumerate}

\subsection{4.4 Monadi: potere compositivo e rischio di deriva}

Le monadi sono dispositivi di composizione potente. In contesti applicativi possono:

\begin{enumerate}
\def\labelenumi{\arabic{enumi}.}
\tightlist
\item
  organizzare pipeline;
\item
  incapsulare effetti;
\item
  semplificare leggi di composizione.
\end{enumerate}

In OCT la domanda è: la composizione iterata resta ordinativamente sana?

Pattern di rischio:

\begin{enumerate}
\def\labelenumi{\arabic{enumi}.}
\tightlist
\item
  ogni passo monadico è localmente accettabile;
\item
  dopo molte iterazioni emerge drift verso B/C;
\item
  la pipeline resta ``ben fatta'' formalmente ma perde funzione reale.
\end{enumerate}

Per questo proponiamo test su iterazioni:

\begin{enumerate}
\def\labelenumi{\arabic{enumi}.}
\tightlist
\item
  \texttt{T\^{}1,\ T\^{}2,\ ...,\ T\^{}k};
\item
  monitoraggio di \texttt{m\_comp} e classe A/B/C;
\item
  criterio di soglia per arresto o revisione.
\end{enumerate}

\subsection{4.5 Comonadi: estrazione di contesto e rischio di rumore}

Le comonadi modellano contesto e coestrazione. In OCT sono preziose quando:

\begin{enumerate}
\def\labelenumi{\arabic{enumi}.}
\tightlist
\item
  il contesto esplicito migliora diagnosi;
\item
  la duplicazione strutturale non esplode in rumore.
\end{enumerate}

Rischio tipico:

\begin{enumerate}
\def\labelenumi{\arabic{enumi}.}
\tightlist
\item
  la duplicazione contestuale cresce più rapidamente della capacità coerentiva;
\item
  il sistema mantiene forma ma perde leggibilità operativa.
\end{enumerate}

Serve quindi una metrica di ``costo contestuale'' da affiancare a \texttt{Coh/Phi}.

\subsection{4.6 Composizione di funtori e preservazione cumulativa}

Come per i morfismi del Capitolo 4, anche per i funtori vale la distinzione:

\begin{enumerate}
\def\labelenumi{\arabic{enumi}.}
\tightlist
\item
  preservazione locale del singolo funtore;
\item
  preservazione globale della catena.
\end{enumerate}

Una catena \texttt{H\_n\ o\ ...\ o\ H\_1} può fallire pur avendo elementi singolarmente buoni, per effetto di interfacce trasformazionali deboli.

Criterio pratico:

\begin{enumerate}
\def\labelenumi{\arabic{enumi}.}
\tightlist
\item
  testare prefissi della catena;
\item
  identificare primo punto di caduta sotto soglia;
\item
  classificare il fallimento come:

  \begin{itemize}
  \tightlist
  \item
    perdita coerentiva;
  \item
    sterilità emergenziale;
  \item
    instabilità contestuale.
  \end{itemize}
\end{enumerate}

\subsection{4.7 Funtori dimenticanti: quando la semplificazione è legittima}

Il funtore dimenticante è spesso indispensabile per ridurre complessità. In OCT non è ``buono'' o ``cattivo'' in astratto:

\begin{enumerate}
\def\labelenumi{\arabic{enumi}.}
\tightlist
\item
  è sicuro su alcune classi;
\item
  è degenerativo su altre.
\end{enumerate}

La regola proposta:

\begin{enumerate}
\def\labelenumi{\arabic{enumi}.}
\tightlist
\item
  dichiarare in anticipo la classe su cui il dimenticante è usato;
\item
  misurare perdita pre/post;
\item
  rifiutare il trasporto forte se la perdita supera soglia critica.
\end{enumerate}

Questo evita che semplificazioni tecniche diventino riduzioni ontologiche implicite.

\subsection{4.8 Naturali tra funtori non equivalenti ordinativamente}

È possibile avere naturali ben definite tra funtori che non sono equivalenti ordinativamente. Questo scenario è importante perché:

\begin{enumerate}
\def\labelenumi{\arabic{enumi}.}
\tightlist
\item
  la presenza di naturali può dare falsa impressione di compatibilità profonda;
\item
  in realtà può esistere divergenza sistemica su robustezza.
\end{enumerate}

Conclusione:

\begin{enumerate}
\def\labelenumi{\arabic{enumi}.}
\tightlist
\item
  naturalezza non è certificato di equivalenza ordinativa.
\end{enumerate}

\subsection{4.9 Regola di accettazione per costrutti trasformazionali}

Proponiamo una regola minima comune a funtori, naturali, aggiunzioni, monadi/comonadi:

un costrutto è ``ordinativamente affidabile'' se:

\begin{enumerate}
\def\labelenumi{\arabic{enumi}.}
\tightlist
\item
  validità classica verificata;
\item
  test ordinativo puntuale superato;
\item
  robustezza su catena o iterazione superata;
\item
  log di fallimenti e limiti pubblicato.
\end{enumerate}

Questa regola unifica formalismo e governance operativa.

\subsection{4.10 Pattern tipici di fallimento nel capitolo}

\begin{enumerate}
\def\labelenumi{\arabic{enumi}.}
\tightlist
\item
  \textbf{F-Fun1}: funtore formalmente corretto ma con drift ordinativo alto.
\item
  \textbf{F-Nat1}: naturale commutativa con perdita emergenziale lungo un lato del quadrato.
\item
  \textbf{F-Adj1}: aggiunzione con ciclo round-trip dissipativo.
\item
  \textbf{F-Mon1}: monade con deriva iterativa verso sterilità.
\item
  \textbf{F-Com1}: comonade con proliferazione contestuale degenerativa.
\item
  \textbf{F-Chain1}: catena trasformazionale con interfaccia debole non individuata.
\end{enumerate}

Questa tipologia favorisce debugging teorico e sperimentale coerente.

\subsection{4.11 Compatibilità con i capitoli precedenti}

Il capitolo è coerente con:

\begin{enumerate}
\def\labelenumi{\arabic{enumi}.}
\tightlist
\item
  Capitolo 4: distinzione locale/globale;
\item
  Capitolo 5: scarto tra equivalenza classica e ordinativa;
\item
  Capitolo 8: spazio A/B/C e invarianti;
\item
  Capitolo 11: criteri di falsificabilità.
\end{enumerate}

Questa compatibilità evita ``isole teoriche'' e mantiene un'unica grammatica del progetto.

\subsection{4.12 Condizione di non-ridondanza del capitolo}

Il capitolo è non ridondante se esiste almeno un caso per ciascuna coppia:

\begin{enumerate}
\def\labelenumi{\arabic{enumi}.}
\tightlist
\item
  funtore classico valido / funtore ordinativamente non affidabile;
\item
  naturale classica valida / naturale ordinativamente fragile;
\item
  aggiunzione classica valida / aggiunzione operativamente dissipativa;
\item
  monade classica valida / monade con deriva iterativa;
\item
  comonade classica valida / comonade con costo contestuale degenerativo.
\end{enumerate}

Il programma D03, D04, D05 e i casi applicativi A01-A03 forniscono la base per soddisfare questa condizione.

\subsection{4.13 Pattern di perdita nelle aggiunzioni}

Le aggiunzioni sono spesso usate per costruire ponti eleganti tra descrizioni diverse dello stesso fenomeno. In OCT è utile distinguere almeno tre pattern di perdita nel ciclo aggiuntivo:

\begin{enumerate}
\def\labelenumi{\arabic{enumi}.}
\item
  \textbf{P-Adj1 (perdita simmetrica lieve)}\\
  andata e ritorno degradano poco e in modo bilanciato; il ciclo resta in area A con margine ridotto.
\item
  \textbf{P-Adj2 (perdita asimmetrica)}\\
  una direzione è robusta, l'altra produce caduta stabile di \texttt{Phi}; il ciclo non è neutro.
\item
  \textbf{P-Adj3 (perdita cumulativa)}\\
  singolo round-trip accettabile, iterazione ripetuta con deriva verso B/C.
\end{enumerate}

Questa classificazione aiuta a evitare un errore comune:

\begin{enumerate}
\def\labelenumi{\arabic{enumi}.}
\tightlist
\item
  valutare l'aggiunzione solo su un passo singolo e non sulla sua dinamica iterata.
\end{enumerate}

\subsection{4.14 Diagnostica duale monade/comonade}

Monadi e comonadi vengono spesso trattate come duali eleganti. In contesti ordinativi la dualità formale non garantisce dualità di stabilità.

Proponiamo una diagnostica duale in quattro controlli:

\begin{enumerate}
\def\labelenumi{\arabic{enumi}.}
\tightlist
\item
  \textbf{controllo di crescita}: come varia il margine su iterazione/coespansione;
\item
  \textbf{controllo di rumore}: quanto costo contestuale viene introdotto;
\item
  \textbf{controllo di recupero}: quanto bene il sistema recupera dopo perturbazione;
\item
  \textbf{controllo di saturazione}: se il sistema converge a punto fisso fertile o sterile.
\end{enumerate}

Esito possibile:

\begin{enumerate}
\def\labelenumi{\arabic{enumi}.}
\tightlist
\item
  monade stabile / comonade fragile;
\item
  comonade stabile / monade derivante;
\item
  entrambe stabili;
\item
  entrambe instabili.
\end{enumerate}

La teoria guadagna precisione perché evita di inferire proprietà operative da sola simmetria sintattica.

\subsection{4.15 Protocollo minimo per validazione trasformazionale}

Per standardizzare la validazione di funtori, naturali, aggiunzioni e (co)monadi, proponiamo un protocollo minimo in sette passi.

\begin{enumerate}
\def\labelenumi{\arabic{enumi}.}
\tightlist
\item
  dichiarazione del dominio e del contesto \texttt{Omega};
\item
  verifica classica del costrutto;
\item
  definizione dei set di test rappresentativi;
\item
  misura di \texttt{Coh/Phi/Delta} su elementi e immagini;
\item
  stress su perturbazioni controllate;
\item
  classificazione dei fallimenti (F-Fun/F-Nat/F-Adj/F-Mon/F-Com/F-Chain);
\item
  decisione finale con stato (\texttt{validated}, \texttt{in\_proof}, \texttt{revise}, \texttt{reject}).
\end{enumerate}

Questo protocollo consente confronto coerente tra gruppi, versioni e domini, riducendo l'ambiguità dei risultati.

\subsection{4.16 Trasporto naturale e conservazione del margine}

Un punto tecnico spesso sottovalutato è la conservazione del margine ordinativo lungo una trasformazione naturale.

Dato \texttt{eta:\ F=\textgreater{}G}, introduciamo:

\texttt{Delta\_m\_eta(X)\ :=\ m\_Omega(GX)\ -\ m\_Omega(FX)}.

Interpretazione:

\begin{enumerate}
\def\labelenumi{\arabic{enumi}.}
\tightlist
\item
  \texttt{Delta\_m\_eta\ \textgreater{}=\ 0} su classe test: trasporto conservativo o migliorativo;
\item
  \texttt{Delta\_m\_eta\ \textless{}\ 0} sistematico: trasporto dissipativo.
\end{enumerate}

Questo indicatore non sostituisce la verifica completa, ma offre una metrica rapida per individuare naturali che sembrano buone formalmente e risultano deboli operativamente.

\subsection{4.17 Composizione di naturali e propagazione di instabilità}

Se \texttt{eta:\ F=\textgreater{}G} e \texttt{theta:\ G=\textgreater{}H}, la composta naturale \texttt{theta\ o\ eta} può introdurre propagazione non lineare dell'instabilità.

Caso tipico:

\begin{enumerate}
\def\labelenumi{\arabic{enumi}.}
\tightlist
\item
  \texttt{eta} induce perdita lieve;
\item
  \texttt{theta} induce perdita lieve;
\item
  la composta supera soglia critica.
\end{enumerate}

Ne segue una regola pratica:

\begin{enumerate}
\def\labelenumi{\arabic{enumi}.}
\tightlist
\item
  testare sempre anche le composte di naturali, non solo i componenti singoli.
\end{enumerate}

\subsection{4.18 Criterio di accettazione forte per catene trasformazionali}

Oltre alla regola minima, definiamo una versione forte:

una catena \texttt{Xi} è \textbf{fortemente affidabile} se:

\begin{enumerate}
\def\labelenumi{\arabic{enumi}.}
\tightlist
\item
  nessun prefisso cade in B/C;
\item
  il margine minimo resta sopra \texttt{eta\_chain\ \textgreater{}\ 0};
\item
  la varianza del margine resta sotto soglia dichiarata;
\item
  i fallimenti osservati sono classificabili e recuperabili senza modifica assiomatica.
\end{enumerate}

Questo criterio è utile in scenari ad alta criticità (pipeline lunghe, sistemi con feedback, applicazioni cross-domain).

\begin{center}\rule{0.5\linewidth}{0.5pt}\end{center}

\section{5. Proposizioni e teoremi locali}

\subsection{Proposizione 6.1 (Necessità del criterio di preservazione ordinativa funtoriale)}

Enunciato:
La sola preservazione classica di composizione/identità non garantisce preservazione ordinativa.

Intuizione:
il funtore può trasportare forma e perdere qualità operativa.

Stato: \texttt{validated}.

\subsection{Proposizione 6.2 (Naturalità ordinativa non automatica)}

Enunciato:
La naturalità classica di \texttt{eta:\ F=\textgreater{}G} non implica naturalità ordinativa.

Intuizione:
quadrati commutativi possono essere ordinativamente asimmetrici per perdita su uno dei percorsi.

Stato: \texttt{in\_review}.

\subsection{Teorema 6.3 (Vincolo di round-trip per aggiunzioni ordinative)}

Ipotesi:

\begin{enumerate}
\def\labelenumi{\arabic{enumi}.}
\tightlist
\item
  \texttt{F\ ⊣\ G} classica;
\item
  test su unità/counità in \texttt{Omega};
\item
  controllo accumulo perdita su iterazioni limitate.
\end{enumerate}

Tesi:
L'aggiunzione è ordinativamente valida solo se il round-trip resta entro margine dichiarato.

Proof sketch:

\begin{enumerate}
\def\labelenumi{\arabic{enumi}.}
\tightlist
\item
  la struttura classica garantisce compatibilità di base;
\item
  il round-trip misura stabilità operativa del trasporto bidirezionale;
\item
  superata la soglia di perdita, l'aggiunzione resta formale ma non affidabile.
\end{enumerate}

Conseguenze:
fornisce criterio operativo per usare aggiunzioni in pipeline reali.

Stato: \texttt{in\_proof}.

\subsection{Teorema 6.4 (Deriva iterativa monadica)}

Ipotesi:

\begin{enumerate}
\def\labelenumi{\arabic{enumi}.}
\tightlist
\item
  monade \texttt{T} classica;
\item
  esiste catena \texttt{T\^{}k} su classe operativa;
\item
  monitoraggio di \texttt{Coh/Phi} per \texttt{k} crescente.
\end{enumerate}

Tesi:
È possibile classificare \texttt{T} come ordinativamente stabile o derivante in base a trend di margine.

Proof sketch:

\begin{enumerate}
\def\labelenumi{\arabic{enumi}.}
\tightlist
\item
  la monade definisce dinamica iterativa naturale;
\item
  il trend del margine compositivo distingue stabilità da deriva;
\item
  la classificazione è indipendente dal singolo passo isolato.
\end{enumerate}

Conseguenze:
introduce controllo longitudinale per monadi in contesti applicativi.

Stato: \texttt{in\_proof}.

\subsection{Proposizione 6.5 (Sicurezza condizionata dei funtori dimenticanti)}

Enunciato:
Un funtore dimenticante può essere sicuro su classi specifiche e degenerativo su altre.

Intuizione:
la perdita informativa non è uniforme; dipende dal tipo di struttura ``dimenticata''.

Stato: \texttt{validated}.

\subsection{Teorema 6.6 (Stabilità di catena trasformazionale)}

Ipotesi:

\begin{enumerate}
\def\labelenumi{\arabic{enumi}.}
\tightlist
\item
  catena \texttt{Xi\ =\ H\_n\ \ o\ \ ...\ \ o\ \ H\_1};
\item
  controllo di preservazione su prefissi;
\item
  assenza di interfacce trasformazionali deboli oltre soglia.
\end{enumerate}

Tesi:
\texttt{Xi} è ordinativamente affidabile sul dominio dichiarato.

Proof sketch:

\begin{enumerate}
\def\labelenumi{\arabic{enumi}.}
\tightlist
\item
  la verifica su prefissi evita masking del fallimento finale;
\item
  l'assenza di interfacce deboli impedisce accumulo nascosto;
\item
  segue affidabilità globale della catena.
\end{enumerate}

Conseguenze:
fissa un criterio pratico per architetture multi-funtoriali.

Stato: \texttt{in\_proof}.

\subsection{Proposizione 6.7 (Non-equivalenza tra dualità formale e dualità di stabilità)}

Enunciato:
La dualità formale tra monade e comonade non implica dualità di comportamento ordinativo.

Intuizione:
la simmetria algebrica può coesistere con asimmetria dei costi contestuali e della deriva.

Stato: \texttt{in\_review}.

\subsection{Teorema 6.8 (Conservazione condizionata del margine naturale)}

Ipotesi:

\begin{enumerate}
\def\labelenumi{\arabic{enumi}.}
\tightlist
\item
  \texttt{eta:\ F=\textgreater{}G} naturale classica;
\item
  esiste classe di test \texttt{K};
\item
  \texttt{Delta\_m\_eta(X)\ \textgreater{}=\ -eps\_m} per ogni \texttt{X\ in\ K}, con \texttt{eps\_m} sotto soglia critica.
\end{enumerate}

Tesi:
Il trasporto naturale \texttt{eta} è ordinativamente accettabile su \texttt{K}.

Proof sketch:

\begin{enumerate}
\def\labelenumi{\arabic{enumi}.}
\tightlist
\item
  la naturalità classica garantisce coerenza diagrammatica;
\item
  il controllo su \texttt{Delta\_m\_eta} limita dissipazione operativa;
\item
  la soglia \texttt{eps\_m} esplicita il margine tollerabile.
\end{enumerate}

Conseguenze:
fornisce un criterio verificabile per approvare naturali in pipeline reali.

Stato: \texttt{in\_proof}.

\begin{center}\rule{0.5\linewidth}{0.5pt}\end{center}

\section{6. Implicazioni operative}

\subsection{6.1 Per progettazione di traduzioni tra modelli}

Ogni traduzione funtoriale dovrebbe esplicitare:

\begin{enumerate}
\def\labelenumi{\arabic{enumi}.}
\tightlist
\item
  cosa preserva formalmente;
\item
  cosa preserva ordinativamente;
\item
  dove introduce perdita controllata.
\end{enumerate}

\subsection{6.2 Per review di trasformazioni naturali}

È utile affiancare al diagramma commutativo una tabella con:

\begin{enumerate}
\def\labelenumi{\arabic{enumi}.}
\tightlist
\item
  classe A/B/C per ciascun percorso;
\item
  differenza di \texttt{Coh/Phi} tra i lati;
\item
  decisione finale (\texttt{accepted}, \texttt{revise}, \texttt{reject}).
\end{enumerate}

\subsection{6.3 Per applicazioni con aggiunzioni}

Nelle applicazioni è consigliato testare round-trip su casi non banali:

\begin{enumerate}
\def\labelenumi{\arabic{enumi}.}
\tightlist
\item
  casi nominali;
\item
  casi al bordo;
\item
  casi rumorosi.
\end{enumerate}

Solo così si evita di validare aggiunzioni ``pulite'' ma non robuste.

\subsection{6.4 Per monadi/comonadi in pipeline iterative}

Serve monitoraggio a finestra:

\begin{enumerate}
\def\labelenumi{\arabic{enumi}.}
\tightlist
\item
  trend su \texttt{k} iterazioni;
\item
  allerta su deriva progressiva;
\item
  trigger automatico di revisione al superamento di soglia.
\end{enumerate}

\subsection{6.5 Per continuità col Capitolo 7}

Le nozioni introdotte qui saranno usate in Capitolo 7 per:

\begin{enumerate}
\def\labelenumi{\arabic{enumi}.}
\tightlist
\item
  strutture monoidali ordinarie e condizionate;
\item
  categorie arricchite con metrica ordinativa;
\item
  fibrations e topos con vincoli di robustezza trasformazionale.
\end{enumerate}

\subsection{6.6 Per governance editoriale}

Ogni release dovrebbe mantenere una sezione ``trasporto trasformazionale'' con:

\begin{enumerate}
\def\labelenumi{\arabic{enumi}.}
\tightlist
\item
  funtori principali usati;
\item
  stato di affidabilità ordinativa;
\item
  fallimenti noti e condizioni di uso.
\end{enumerate}

Questo riduce regressioni e migliora leggibilità del manoscritto evolutivo.

\subsection{6.7 Per integrazione con benchmark futuri}

Nel disegno dei benchmark dei Capitoli 12-13 è utile introdurre suite dedicate:

\begin{enumerate}
\def\labelenumi{\arabic{enumi}.}
\tightlist
\item
  \texttt{FUN\_SUITE}: preservazione funtoriale;
\item
  \texttt{NAT\_SUITE}: coerenza dei quadrati naturali con margine;
\item
  \texttt{ADJ\_SUITE}: round-trip e perdita cumulativa;
\item
  \texttt{MONCOM\_SUITE}: deriva iterativa e costo contestuale.
\end{enumerate}

Questa segmentazione rende più semplice attribuire cause e responsabilità dei fallimenti.

\subsection{6.8 Per comunicazione dei risultati}

Per evitare ambiguità, ogni risultato trasformazionale dovrebbe riportare:

\begin{enumerate}
\def\labelenumi{\arabic{enumi}.}
\tightlist
\item
  costrutto classico coinvolto;
\item
  classe di test usata;
\item
  indicatori di preservazione (\texttt{DeltaCoh}, \texttt{DeltaPhi}, \texttt{Drift}, \texttt{Delta\_m});
\item
  stato finale e motivazione.
\end{enumerate}

Questo formato consente lettura rapida da parte di matematici, ingegneri e revisori interdisciplinari.

\begin{center}\rule{0.5\linewidth}{0.5pt}\end{center}

\section{7. Limiti e non-claim}

Limiti espliciti:

\begin{enumerate}
\def\labelenumi{\arabic{enumi}.}
\tightlist
\item
  il capitolo non chiude tutte le prove complete su aggiunzioni e monadi/comonadi;
\item
  la metrica di costo contestuale per comonadi è introdotta in forma minima;
\item
  la stabilità di catene lunghe richiede campagne di benchmark dedicate.
\end{enumerate}

Failure mode riconosciuti:

\begin{enumerate}
\def\labelenumi{\arabic{enumi}.}
\tightlist
\item
  confondere preservazione classica con preservazione ordinativa;
\item
  validare naturalità senza analisi dei due percorsi;
\item
  usare monadi iterative senza controllo di deriva;
\item
  assumere sicurezza dei funtori dimenticanti fuori dalla classe validata;
\item
  estendere risultati di un dominio a un altro senza mapping contestuale.
\end{enumerate}

Box C - Non-claim

Questo capitolo non implica:

\begin{enumerate}
\def\labelenumi{\arabic{enumi}.}
\tightlist
\item
  che ogni funtore classico debba diventare ordinativamente ammissibile;
\item
  che ogni aggiunzione utile teoricamente sia automaticamente utile operativamente;
\item
  che monadi/comonadi classiche siano sempre stabili in scenari reali iterativi.
\end{enumerate}

\begin{center}\rule{0.5\linewidth}{0.5pt}\end{center}

\section{8. Claim Ledger (S0-S3)}

{\def\LTcaptype{none} % do not increment counter
\begin{longtable}[]{@{}lllll@{}}
\toprule\noalign{}
ID & Claim & Tipo & Evidenza & Stato \\
\midrule\noalign{}
\endhead
\bottomrule\noalign{}
\endlastfoot
C6.1 & La preservazione funtoriale classica è necessaria ma non sufficiente per affidabilità ordinativa & formale & S1 & validated \\
C6.2 & Naturalità classica e naturalità ordinativa possono divergere & formale & S2 & in\_review \\
C6.3 & Le aggiunzioni richiedono test di round-trip ordinativo per uso forte & formale-operativo & S2 & in\_proof \\
C6.4 & Le monadi richiedono analisi di deriva su iterazione, non solo verifica locale & formale-operativo & S2 & in\_proof \\
C6.5 & I funtori dimenticanti sono sicuri solo su classi dichiarate & metodologico & S1 & validated \\
C6.6 & La stabilità di catene trasformazionali richiede controllo su prefissi e interfacce & formale-operativo & S2 & in\_proof \\
C6.7 & Il capitolo fornisce il ponte tecnico necessario verso monoidali/enriched/fibrations/topos ordinativi & programmatico & S1 & in\_review \\
\end{longtable}
}

\begin{center}\rule{0.5\linewidth}{0.5pt}\end{center}

\section{9. Bridge al Capitolo 7}

Il Capitolo 6 ha esteso il livello trasformazionale della teoria:

\begin{enumerate}
\def\labelenumi{\arabic{enumi}.}
\tightlist
\item
  funtori e naturali con vincolo ordinativo;
\item
  aggiunzioni con controllo di round-trip;
\item
  monadi/comonadi con monitoraggio di deriva e costo contestuale.
\end{enumerate}

Il Capitolo 7 porterà questa estensione a strutture categoriali avanzate:

\begin{enumerate}
\def\labelenumi{\arabic{enumi}.}
\tightlist
\item
  monoidali;
\item
  arricchite;
\item
  fibrations;
\item
  topos e dinamiche categoriali.
\end{enumerate}

In sintesi: il Capitolo 6 stabilisce come trasportare la validità ordinativa; il Capitolo 7 mostrerà come farla vivere in strutture ad alta complessità.

\begin{center}\rule{0.5\linewidth}{0.5pt}\end{center}

\section{Riferimenti}

Riferimenti di framework interno:

\begin{enumerate}
\def\labelenumi{\arabic{enumi}.}
\tightlist
\item
  Ghioni, F. (2026). \texttt{TE\_CORE\_v5.1.md}.
\item
  Ghioni, F. (2026). \texttt{Ordinative\_Set\_Theory\_OST\_A\_Concise\_Guide\_For\_AI\_v2\_1.md}.
\item
  \texttt{OCT\_FOUNDATIONAL\_BOOK\_CHAPTER\_04\_v1\_0.md}.
\item
  \texttt{OCT\_FOUNDATIONAL\_BOOK\_CHAPTER\_05\_v1\_0.md}.
\item
  \texttt{OCT\_FOUNDATIONAL\_BOOK\_CHAPTER\_08\_v1\_0.md}.
\item
  \texttt{OCT\_FOUNDATIONAL\_BOOK\_CHAPTER\_09\_v1\_0.md}.
\end{enumerate}

Riferimenti primari:

\begin{enumerate}
\def\labelenumi{\arabic{enumi}.}
\tightlist
\item
  Mac Lane, S. (1998). \emph{Categories for the Working Mathematician}.
\item
  Riehl, E. (2016). \emph{Category Theory in Context}.
\item
  Spivak, D. (2014). \emph{Category Theory for the Sciences}.
\end{enumerate}
